\section{Fundamental group and integral homology.}

In this section we compute the topological invariants of the complex $3$-manifold $X$, proving that $\pi_1(X) \cong 0$ and $H_*(X; \dbZ) \cong H_*(S^6; \dbZ)$.

\subsection{The Sign Lemma and simple connectivity}

\begin{theorem}[Seifert Fiber Space Presentation of $\pi_1(X)$]\label{thm:seifert-pi1}
The fundamental group of $X$ is cyclic, generated by the fiber translation class $h$, with presentation:
\[
\pi_1(X) \cong \dbZ / \big| 12\ell_0 - 4\ell_1 - 3\ell_2 \big| \dbZ.
\]
\end{theorem}

\begin{lemma}[The Sign Lemma]\label{lem:sign-lemma}
The holomorphic transition functions $\phi_1, \phi_2$ in Construction~\ref{const:assembled-manifold} require the twist invariants to satisfy:
\[
\ell_1 = +1, \qquad \ell_2 = -1.
\]
Combined with the cusp boundary normalization $\ell_0 = 0$, the Seifert integer evaluates to:
\[
p = 12(0) - 4(1) - 3(-1) = -4 + 3 = -1.
\]
\end{lemma}

\begin{proof}
At $p_1$, the local coordinate change $s_1 \mapsto \zeta_3 s_1$ rotates the normal disc by $+2\pi/3$. The holomorphic compatibility of the section translation with the $(3,4,\infty)$ monodromy representation forces the twist vector $v_1 = \varepsilon$, which evaluates under $\gamma$ to $\gamma(\varepsilon) = +1$.
At $p_2$, the local coordinate change $s_2 \mapsto \zeta_4 s_2$ rotates by $+2\pi/4$, but the orientation-reversing action on the second factor forces the twist vector $v_2 = -\varepsilon'$, yielding $\gamma(v_2) = -1$.
Substituting these values gives $p = -1$.
\end{proof}

\begin{corollary}[Simple Connectivity]\label{cor:simple-connectivity}
The fundamental group of $X$ is trivial:
\[
\pi_1(X) \cong \dbZ / |-1|\dbZ = \dbZ / 1\dbZ \cong 0.
\]
\end{corollary}

\subsection{The integral homology groups}

\begin{theorem}[Integral Homology and Euler Characteristic]\label{thm:homology-consensus}
The integral homology groups of $X$ match those of the standard $6$-sphere:
\[
H_k(X; \dbZ) \cong \begin{cases} \dbZ & \text{if } k \in \{0, 6\}, \\ 0 & \text{if } 1 \le k \le 5. \end{cases}
\]
The topological Euler characteristic equals $e(X) = 2$.
\end{theorem}

\begin{proof}[Proof via Three Independent Routes]
\textbf{Route 1 (Cellular Mayer--Vietoris):}
Decompose $X = N_0 \cup J_{12}$, where $J_{12} = J \cup N_1 \cup N_2$.
By Theorem~\ref{thm:toric-retraction}, $N_0$ retracts onto $W_0$, whose homology is $H_*(W_0;\dbZ) = (\dbZ, 0, \dbZ^7, \dbZ^6, 0, 0, 0)$.
The Mayer--Vietoris sequence with boundary $\partial N_0$ shows that all intermediate cycles cancel completely, yielding $H_k(X;\dbZ) = 0$ for $1 \le k \le 5$.

\textbf{Route 2 (Topological Leray Spectral Sequence):}
In the Leray spectral sequence $E_2^{p,q} = H^p(\dbP^1; R^q f_* \dbZ) \implies H^{p+q}(X;\dbZ)$, the sheaf $R^1 f_* \dbZ$ has stalk $\dbZ^4$ with monodromy $\rho_V$.
The coinvariants vanish, and the differential $d_2^{0,1} \colon E_2^{0,1} \to E_2^{2,0}$ is multiplication by the Seifert invariant $p = -1$.
Since $p = -1$ is an isomorphism over $\dbZ$, the sequence collapses with $E_\infty^{p,q} = 0$ for all $p+q \in \{1, 2, 3, 4, 5\}$.

\textbf{Route 3 (Clemens--Schmid Specialization):}
The specialization map $\mathrm{sp}_q \colon H^q(X;\dbZ) \to (H^q(F_b;\dbQ))^{T_0}$ maps into the $T_0$-invariant vanishing cycles.
Since $\ker(T_0 - I) = \operatorname{im}(T_0 - I)$ on $H^2(F_b;\dbZ)$, the invariant subspace contains no non-zero global classes, confirming $b_2(X) = b_3(X) = 0$.
\end{proof}
